% This supplementary material provides useful definitions and propositions, necessary lemmas with detailed proofs, and the complete derivations for the main theoretical results and parameters presented in the primary work.
\section{Preliminaries}
\label{Appendix: Preliminaries and notations}
% The proximal mapping admits the following
% properties.
% \begin{proposition}[Properties of proximal operators
% \cite{Rockafellar1998VariationalA}]
%     Let $f : \mathbb{R}^{q}\rightarrow \mathbb{R}\cup \{+\infty\}$ be proper and
%     closed. Then, the following holds.
%     \begin{enumerate}[label=(\roman{*})]
%         \item The set $\operatorname{prox}_{\tau f}(\bx)$ is nonempty and compact
%             for any $\bx \in \mathbb{R}^{q}$ and $\tau > 0$.

%         \item If $f$ is convex, then $\operatorname{prox}_{\tau f}(\bx)$
%             contains exactly one value for any $\bx \in \mathbb{R}^{q}$ and
%             $\tau > 0$.
%     \end{enumerate}
% \end{proposition}

We introduce the following definitions related to subdifferentials.
\begin{definition}[Subdifferentials \cite{Rockafellar1998VariationalA}]
    \label{definition: subgradient} Let $f : \mathbb{R}^{q}\to (-\infty, +\infty
    ]$ be a proper closed function.
    \begin{enumerate}[label=(\roman{*})]
        \item For a given $\bx \in \operatorname{dom}f$, the Fréchet
            subdifferential, or simply of $f$ at $\bx$, written
            $\hat{\partial}f(\bx)$, is the set of all vectors $\bu \in \mathbb{R}
            ^{q}$ that satisfy
            \[
                \liminf_{\by \neq \bx, \by \to \bx}\frac{f(\by) - f(\bx) -
                \langle \bu, \by - \bx \rangle}{\|\by - \bx\|}\geq 0.
            \]

        \item The limiting-subdifferential, or simply the subdifferential, of $f$
            at $\bx$, written $\partial f(\bx)$, is defined through the following
            closure process
\[
\partial f(\bx) := \{ \bu \in \mathbb{R}^{q}: \exists \bx^{k}\to \bx, f(\bx^{k}) \to f(\bx), \bu^{k}\in \hat{\partial}f(\bx^{k}) \to \bu \text{ as }k \to \infty \}.
\]
    \end{enumerate}
\end{definition}
Notionally, $\hat{\partial}f(\bx) = \partial f(\bx) = \emptyset$ for all $\bx \notin
\operatorname{dom}(f)$. From definition \cref{definition: subgradient} it
follows that $\hat{\partial}f(\bx) \subset \partial f(\bx)$ for all
$\bx \in \operatorname{dom}(f)$. Moreover, $\hat{\partial}f(\bx)$ is closed
convex while $\partial f(\bx)$ is merely closed. In particular, if $f$ is convex,
then
\[
    \hat{\partial}f(\bx) = \partial f(\bx) = \{ \bd \in \mathbb{R}^{q}\mid f(\by
    ) \geq f(x) + \langle \bd, \by - \bx \rangle ,\,\forall \by \in \mathbb{R}
    ^{q}\}.
\]
The following proposition lists several useful properties of subgradients.
\begin{proposition}[Properties of subdifferential\cite{Rockafellar1998VariationalA}]
    Let $f : \mathbb{R}^{q_1}\rightarrow (-\infty, \infty]$ and $g : \mathbb{R}^{q_2}
    \rightarrow (-\infty, \infty]$ be proper and lower semicontinuous, and let
    $h : \mathbb{R}^{q_1}\rightarrow \mathbb{R}$ be continuously differentiable.
    Then the following results hold:
    \begin{enumerate}[label=(\roman{*})]
        \item For any $\bx \in \mathbb{R}^{q_1}$, we have
            $\hat{\partial}f(\bx) \subseteq \partial f(\bx)$. If, in addition, $f$
            is convex, we have $\hat{\partial}f(\bx) = \partial f(\bx)$.

        \item For $\bx \in \mathbb{R}^{q_1}$ with $f(\bx) < \infty$, it holds
           \[
\begin{aligned}
    \hat{\partial}(f + h)(\bx) &= \hat{\partial}f(\bx) + \nabla h(\bx) \quad \text{and}\\
    \partial (f + h)(\bx) &= \partial f(\bx) + \nabla h(\bx).
\end{aligned}
\]
        \item If $\sigma(\bx_{1}, \bx_{2}) = f_{1}(\bx_{1}) + f_{2}(\bx_{2})$,
            then
            \[
                \begin{pmatrix}
                    \hat{\partial}_{\bx_1}f_{1}(\bar{\bx}_{1}) \\
                    \hat{\partial}_{\bx_2}f_{2}(\bar{\bx}_{2})
                \end{pmatrix}
                \subseteq \hat{\partial}\sigma(\bar{\bx}_{1}, \bar{\bx}_{2}) \quad
                \text{and}\quad
                \begin{pmatrix}
                    \partial_{\bx_1}f_{1}(\bar{\bx}_{1}) \\
                    \partial_{\bx_2}f_{2}(\bar{\bx}_{2})
                \end{pmatrix}
                \subseteq \partial \sigma(\bar{\bx}_{1}, \bar{\bx}_{2}).
            \]
    \end{enumerate}
\end{proposition}
For \Cref{alg:our_algorithm_sgdm}, the key relationships among iterates via the momentum acceleration framework in \eqref{Eq: the update of x in SGDM} are established as follows.
\begin{proposition}[Lemma 5.2 in \cite{InertialProximalAlternating}]
    \label{Pp: relations of xyz} Let $\{\bx^{k}\}_{k \in \mathbb{N}}$,
    $\{\by^{k}\}_{k \in \mathbb{N}}$, and $\{\bz^{k}\}_{k \in \mathbb{N}}$ be sequences
    generated by \eqref{Eq: the update of x in SGDM}. For all $k \in \mathbb{N}$,
    we have:
    \begin{enumerate}[label=(\roman{*})]
        \item $\|\bx^{k}- \by^{k}\|^{2}= \alpha_{k}^{2}\|\bx^{k}- \bx^{k-1}\|^{2}$,

        \item $\|\bx^{k}- \bz^{k}\|^{2}= \beta_{k}^{2}\|\bx^{k}- \bx^{k-1}\|^{2}$,

        \item $\|\bx^{k+1}- \by^{k}\|^{2}\geq 2(1 - \alpha_{k})\|\bx^{k+1}- \bx^{k}
            \|^{2}+ 2\alpha_{k}(\alpha_{k}- 1)\|\bx^{k}- \bx^{k-1}\|^{2}$.
    \end{enumerate}
\end{proposition}
% For  \Cref{alg:our_algorithm_sgdm},  we establish key relationships among iterates via the momentum acceleration framework in \eqref{Eq: the update of x in SGDM}.
% \begin{proposition}[Lemma 5.2 in \cite{InertialProximalAlternating}]
%     \label{Pp: relations of xyz} Let $\{\bx^{k}\}_{k \in \mathbb{N}}$,
%     $\{\by^{k}\}_{k \in \mathbb{N}}$, and $\{\bz^{k}\}_{k \in \mathbb{N}}$ be sequences
%     generated by \eqref{Eq: the update of x in SGDM}. For all $k \in \mathbb{N}$,
%     we have:
%     \begin{enumerate}[label=(\roman{*})]
%         \item $\|\bx^{k}- \by^{k}\|^{2}= \alpha_{k}^{2}\|\bx^{k}- \bx^{k-1}\|^{2}$,

%         \item $\|\bx^{k}- \bz^{k}\|^{2}= \beta_{k}^{2}\|\bx^{k}- \bx^{k-1}\|^{2}$,

%         \item $\|\bx^{k+1}- \by^{k}\|^{2}\geq 2(1 - \alpha_{k})\|\bx^{k+1}- \bx^{k}
%             \|^{2}+ 2\alpha_{k}(\alpha_{k}- 1)\|\bx^{k}- \bx^{k-1}\|^{2}$.
%     \end{enumerate}
% \end{proposition}
% For \Cref{alg:our_algorithm_storm}, we employ the STORM (STochastic Recursive Momentum)  variance reduction estimator \cite{cutkoskyMomentumBasedVarianceReduction2019}.
% In stochastic optimization problems where $F(\bx) = \mathbb{E}_{\xi}[F(\bx;\xi)]$ with $\xi \in \Xi$, standard momentum methods use
% \[
%     \bd^{k}= \rho_{k}\nabla F(\bx^{k}, \xi^{k}) + (1 - \rho_{k}) \bd^{k-1}, \text{
%     where }\rho_{k}\in [0, 1],
% \] yielding a biased estimate of $\nabla F(\bx^{k})$. 
% STORM adds a
% correction term, $(1 - \rho_{k}) (\nabla F(\bx^{k}, \xi^{k}) - \nabla F(\bx^{k-1}
% , \xi^{k})$, resulting in:
% \[
%     \begin{aligned}
%          & \bd^{k}=\left\{\begin{array}{ll}\nabla F\left(\bx^{k}, \xi^{k}\right), & \text{if } k=1, \\ \nabla F\left(\bx^{k}, \xi^{k}\right) + \left(1-\rho_k\right)\left( \bd^{k-1} - \nabla F\left(\bx^{k-1}, \xi^{k}\right) \right), & \text{if } k>1,\end{array}\right. \\
%          & \bx^{k+1}=\bx^{k}-\eta_{k}\bd^{k}.
%     \end{aligned}
% \]
% This exploits the smoothness of $F(\cdot, \xi)$, leading to efficient
% variance reduction. When $\rho_{k}= 1$, $\bd^{k}$ reduces to vanilla mini-batch SGD. We set $\rho_{k}\in (0,1)$ for $k \geq 1$ in \Cref{alg:our_algorithm_storm}.
% \newpage
\section{Discussions on \Cref{Ass: main assumptions}(ii) in \Cref{Section: Convergence analysis}}

\label{Appendix: discussion on assumptions} 
This section verifies that the smoothness constant $L_F$ in \Cref{Ass: main assumptions}(ii) is independent of the sample size $n$, which is essential for our convergence analysis.

For the smooth properties of mappings in our problem, we consider the nonlinear mapping
$\mathcal{D}(\bx) = (\bD_{1}(\bx), \ldots, \bD_{n}(\bx))$ from $\mathcal{X}\subseteq
\RR^{d_{\bx}}$ to $\RR^{n \times d}$, and make the following assumption:

\begin{assumption}
    [Properties of $\mathcal{D}$] \label{Ass: the properties of D} Suppose there
    exist constants $L_{\mathcal{D}},\, L_{\nabla \mathcal{D}}>0$ such that for any
    $\bx, \bx' \in \mathcal{X}$:
    $\|\mathcal{D}(\bx) - \mathcal{D}(\bx')\| \leq L_{\mathcal{D}}\|\bx - \bx'\|$
    and
    $\|\nabla \mathcal{D}(\bx) - \nabla \mathcal{D}(\bx')\| \leq L_{\nabla
    \mathcal{D}}\|\bx - \bx'\|$.
\end{assumption}

Combining \Cref{Ass: main assumptions}(i) and \Cref{Ass: the properties of D},
we can establish that there exist constants
$C_{\mathcal{D}}, G_{\mathcal{D}}, C_{\mathcal{D}B}> 0$ such that:
\[
    \|\mathcal{D}(\bx)\| \leq C_{\mathcal{D}}, \, \|\nabla \mathcal{D}(\bx)\|
    \leq G_{\mathcal{D}}, \, \text{and}\quad \|\mathcal{D}(\bx) - \bB\| \leq C
    _{\mathcal{D}B}.
\]
Based on these assumptions, we can establish the smoothness properties of the composite
function $F$. First, we observe that $P(\bx, \bB)$ is $\frac{\gamma}{n}$-smooth with
respect to $\bB$ and $\frac{\gamma}{n}(L_{\nabla \mathcal{D}}(C_{\mathcal{D}}+ C_{\bB}
) + G_{\mathcal{D}}L_{\mathcal{D}})$-smooth with respect to $\bx$. Combining
these properties, $P(\bx, \bB)$ has an overall smoothness constant:
\[
    L_{s}= \frac{\gamma}{n}\sqrt{2\max\{L^{2}_{\nabla \mathcal{D}}C^{2}_{\mathcal{D}B}+
    (2G_{\mathcal{D}}^{2}+ 1)L^{2}_{\mathcal{D}}, 2G_{\mathcal{D}}^{2}+ 1\}}= \mathcal{O}
    (1)
\]

Consequently, $F(\bx, \bB)$ is $L_{F}$-smooth with respect to $(\bx, \bB)$, where
$L_{F}= \mathcal{O}(1)$. An important observation is that, despite $C_{\mathcal{D}}$,
$G_{\mathcal{D}}$ and $C_{\mathcal{D}B}$ potentially being
$\mathcal{O}(\sqrt{n})$, the overall smoothness constant $L_{F}$ remains independent
of the sample size $n$, which is crucial for our convergence analysis.
% This smoothness property implies the following descent lemma for $F$ with
% respect to $\bx$ and $\bB$:
% \begin{align}
% \label{Eq: F smoothness with x}
%     F(\bx', \bB) \leq F(\bx, \bB) + \langle \nabla_{\bx}F(\bx, \bB), \bx' -\bx\rangle
%     + \frac{L_{F}}{2}\| \bx - \bx'\|^{2}
% \end{align}
% for any $\bB \in \calB$; and similarly for $\bB$:
% {\small 
% \begin{align}
% \label{Eq: F smoothness with B}
%     F(\bx, \bB') \leq F(\bx, \bB) + \langle \nabla_{\bB}F(\bx, \bB), \bB' -\bB\rangle
%     + \frac{L_{F}}{2}\| \bB - \bB'\|^{2}
% \end{align}} for any $\bx \in \calX$.

\section{Complexity analysis of DualHash-StoM (\texorpdfstring{\Cref{alg:our_algorithm_sgdm}}{Algorithm 1})}
\label{Appendix: the complexity analysis of DualHash-StoM}
We adopt the notation and intermediate results from \Cref{Section: Convergence analysis}.
Let us introduce the notation that will be used throughout our analysis:
\begin{align*}
    \bxi^{k}                & = \{\xi_{j}^{k}, j \in \calJ_{k}\},        & \text{(mini-batch at iteration $k$)}                 \\
    \bxi^{[k]}              & = \{\bxi^{1}, \bxi^{2}, \ldots, \bxi^{k}\}, & \text{(history of mini-batches up to iteration $k$)} \\
    \EE_{\bxi^{[k]}}[\cdot] & = \EE [ \cdot | \bxi^{[k]}] .              & \text{(conditional expectation given history)}
\end{align*}

For $\bd^{k}\in \partial \calL(\bx^{k}, \bB^{k}, \bm{\Lambda}^{k})$, let us
denote $\bd^{k}= (\bd_{\bx}^{k}, \bd_{\bB}^{k}, \bd_{\bm{\Lambda}}^{k})$ and
\begin{align}
    % \label{Eq: the components of the subgradient}
    \begin{cases}&\bd_{\bx}^{k}= \nabla_{\bx}F(\bx^{k}, \bB^{k}),\\&\bd_{\bB}^{k}= \nabla_{\bB}F(\bx^{k}, \bB^{k}) + \bm{\Lambda}^{k}, \\&\bd_{\bm{\Lambda}}^{k}= \bB^{k}- \calG_{\bm{\Lambda}}^{k},\end{cases} \notag
\end{align}
where $\calG_{\bm{\Lambda}}^{k}\in \partial h^{*}(\bm{\Lambda}^{k})$.
We first present a lemma concerning a property of a continuously differentiable function with a Lipschitz continuous gradient, followed by a proposition that details the interrelations among the sequences $\{\bx^{k}\}_{k \in \NN},\, \{\bm{y}^{k}\}_{k \in \NN},\, \{\bm{z}^{k}\}_{k \in \NN}$ generated by the algorithm.
% In particular, $\Delta_{\bx}^{k}$ and $\Delta_{\bB}^{k}$ are as defined therein, and we refer to \Cref{Lemma: dual boundedness} and \Cref{Lemma: One-iteration progress of B and Lambda} for the bounds on dual variable dynamics and one-iteration Lagrangian progress on $\bB$ and $\bLmbd$. The gradient error is defined as $\be^{k}= \bG^{k}- \nabla_{\bx}F(\bz^{k}, \bB^{k})$, satisfying $\EE_{\bxi^{[k]}}[\|\be^{k}\|^{2}] \leq \sigma^{2}/b_{k}$ under \Cref{Ass: main assumptions}~(iii).
\begin{lemma}
    \label{Lemma: auxiliary lemma for f} Let $f(\bx): \mathbb{R}^{q}\rightarrow \mathbb{R}$
    be a continuously differentiable function with $L$-lipschitz continuous gradient.
    Then, for any $\bx,\, \by,\,\bz \in \mathbb{R}^{q}$ and any $\bx^{+}\in \mathbb{R}
    ^{q}$ defined by
    \[
        \bx^{+}= \by - \eta \nabla f(\bz), \hspace{5pt}\eta > 0,
    \]it holds that
    \begin{align}
        \label{Eq: auxiliary lemma for f}f(\bx^{+}) \leq f(\bx) + \langle \nabla f(\bx) - \nabla f(\bz), \bx^{+}- \bx \rangle + \frac{L}{2}\|\bx^{+}- \bx\|^{2}+ \frac{1}{2\eta}\|\bx - \by\|^{2}- \frac{1}{2\eta}\|\bx^{+}- \by\|^{2}.
    \end{align}
\end{lemma}
\begin{proof}
    By the definition of $\bx^{+}$, we find that
    \[
        \bx^{+}= \argmin\limits_{\bu \in \mathbb{R}^d}\{ \langle \nabla f(\bz), \bu
        - \by \rangle + \frac{1}{2\eta}\| \bu - \by\|^{2}\},
    \]
    and hence, in particular, by taking $\bu = \bx$, we obtain
    \begin{align*}
        0 \leq \langle \nabla f(\bz), \bx - \bx^{+}\rangle + \frac{\tau}{2}\| \bx - \by\|^{2}- \frac{1}{2\eta}\| \bx^{+}- \by\|^{2}.
    \end{align*}
    Invoking the descent lemma for a smooth function yields the result. 
\end{proof}
\subsection{Proof of \Cref{Lemma: One-iteration progress of x in SGDM}}
\label{Subsubsec: proof x progress SGDM}

\begin{proof}
    Recalling the equivalent relation:
    \begin{align}
        \label{Eq: the update of x one progress in StoM}\mathcal{L}(\bx^{k+1}, \bm{B}^{k}, \bm{\Lambda}^{k}) - \mathcal{L}(\bx^{k}, \bm{B}^{k}, \bm{\Lambda}^{k}) = F(\bx^{k+1}, \bm{B}^{k}) - F(\bx^{k}, \bm{B}^{k}).
    \end{align}
    Then, applying the smoothness of $F$ with respect to $\bx$ and \eqref{Eq: auxiliary lemma for f}
    yields
    \begin{align}
        \label{Eq: the update of theta core}F(\bx^{k+1}, \bm{B}^{k}) \leq & F(\bx^{k}, \bm{B}^{k}) + \langle \nabla_{\bx}F(\bx^{k}, \bm{B}^{k}) - \bG^{k}, \bx^{k+1}- \bx^{k}\rangle \notag                              \\
        +                                                                 & \frac{L_{F}}{2}\|\bx^{k+1}- \bx^{k}\|^{2}+ \frac{1}{2\eta_{k}}\|\bx^{k}- \bm{y}^{k}\|^{2}- \frac{1}{2\eta_{k}}\|\bx^{k+1}- \bm{y}^{k}\|^{2}.
    \end{align}
    For the inner product term in \eqref{Eq: the update of theta core}, using
    the facts that
    $\langle a, b \rangle \leq \frac{s}{2}\|a\|^{2}+ \frac{1}{2s}\|b\|^{2}$ for
    any $s > 0$ and $\| a - c \|^{2}\leq 2\|a - b \|^{2}+ 2\|b -c\|^{2}$ yields
    { \begin{align}\label{Eq: Inner product for theta}&\langle \nabla_{\bx}F(\bx^{k}, \bm{B}^{k}) - \bG^{k}, \bx^{k+1}- \bx^{k}\rangle \notag \\ \leq&\,\frac{s}{2}\|\bx^{k+1}- \bx^{k}\|^{2}+ \frac{1}{s}\|\nabla_{\bx}F(\bx^{k}, \bm{B}^{k}) - \nabla_{\bx}F(\bm{z}^{k}, \bm{B}^{k}) \|^{2}+ \frac{1}{s}\|\nabla_{\bx}F(\bm{z}^{k}, \bm{B}^{k}) - \bG^{k}\|^{2}\notag \\ \leq&\, \frac{s}{2}\|\bx^{k+1}- \bx^{k}\|^{2}+ \frac{L^{2}_{F}}{s}\|\bx^{k}- \bm{z}^{k}\|^{2}+ \frac{1}{s}\|\nabla_{\bx}F(\bm{z}^{k}, \bm{B}^{k}) - \bG^{k}\|^{2}.\end{align}}
    Setting $s = L_{F}$ in \eqref{Eq: Inner product for theta} and substituting into
    \eqref{Eq: the update of theta core}, then applying \cref{Pp: relations of xyz},
    we obtain
    \begin{align}
        F(\bx^{k+1}, \bm{B}^{k}) & \leq F(\bx^{k}, \bm{B}^{k}) - (\frac{1 - \alpha_{k}}{\eta_{k}}- 2L_{F})\frac{\|\bx^{k+1}- \bx^{k}\|^{2}}{2}\notag                                                                                           \\
                                 & +\,(\frac{\alpha_{k}}{\eta_{k}}+ 2L_{F}\beta_{k}^{2})\frac{\|\bx^{k}- \bx^{k-1}\|^{2}}{2}+ \frac{1}{L_{F}}\|\nabla_{\bx}F(\bm{z}^{k}, \bm{B}^{k}) - \bG^{k}\|^{2}. \label{Eq: the update of x of L in StoM}
    \end{align}
    Substituting \eqref{Eq: the update of x of L in StoM} into \eqref{Eq: the update of x one progress in StoM}
    yields \eqref{Eq: One-iteration progress of x in SGDM} and completes the proof.
    
    % \begin{align}
    %     \label{item: x}
    %     &\mathcal{L}(\bx^{k+1}, \bm{B}^k, \bm{\Lambda}^k) -\mathcal{L}(\bx^k, \bm{B}^k, \bm{\Lambda}^k)  \notag \\
    %     \leq &\,-(\frac{1 - \alpha_k}{\eta_k}- 2L_F) \frac{ \|\bx^{k+1} - \bx^k \|^2}{2}
    %     + (\frac{\alpha_k}{\eta_k}+ 2L_F\beta_k^2) \frac{\|\bx^k - \bx^{k-1} \|^2}{2} \notag \\
    %     \quad&+ \frac{1}{L_F} \|\nabla_{\bx} F(\bm{z}^k, \bm{B}^k) - \bG^k \|^2.
    % \end{align} 
\end{proof}
\subsection{Proof of \Cref{Lemma: approximate sufficient descent property in SGDM}}
\label{Subsubsec: proof descent SGDM}

\begin{proof}
    Substituting one-iteration progress \eqref{Eq: One-iteration progress of x in SGDM}
    and \eqref{item: B and Lambda} into \eqref{Eq: One-iteration progress decomposition}
    , we obtain
    \begin{align}
        \label{Eq: One-iteration progress in sgdm} & \calL (\bx^{k+1}, \bB^{k+1}, \bLmbd^{k+1}) - \calL (\bx^{k}, \bB^{k}, \bLmbd^{k}) \notag                                                                                                   \\
        \leq \,                                    & K_{1}\Delta_{\bB}^{k+2}- K_{2}\Delta_{\bB}^{k+1}+ K_{3}\Delta_{\bB}^{k}+ K_{4}\Delta_{\bx}^{k+2}- (K_{5}- K_{3})\Delta_{\bx}^{k+1}+ K_{6}\Delta_{\bx}^{k}+ \frac{1}{L_{F}}\|\be^{k}\|^{2}.
    \end{align}
    Applying the definition of $\Psi_m$ from \eqref{Eq: the Lyapunov function in SGDM}
    to \eqref{Eq: One-iteration progress in sgdm}, rearranging, and taking the
    conditional expectation on both sides, we obtain
    \begin{align}
        \label{the descent property of Lyapunov func in SGDM first} & \EE_{\bxi^{[k+1]}}\left[\Psi_m(\mathcal{Q}^{k+1})\right] - \EE_{\bxi^{[k]}}\left[\Psi_m(\mathcal{Q}^{k})\right]\notag                                                                               \\
        \leq                                                        & - (C_{1}- K_{1}) \EE_{\bxi^{[k+1]}}[\Delta_{\bm{B}}^{k+2}] -(K_{2}- C_{1}- C_{2}) \EE_{\bxi^{[k+1]}}[\Delta_{\bm{B}}^{k+1}] - (C_{2}- K_{3}) \EE_{\bxi^{[k]}}[\Delta_{\bm{B}}^{k}] \notag                     \\
                                                                    & - (C_{3}-K_{4}) \EE_{\bxi^{[k+1]}}[\Delta_{\bx}^{k+2}]- (K_{5}-K_{3}-C_{3}-C_{4}) \EE_{\bxi^{[k+1]}}[\Delta_{\bx}^{k+1}] - (C_{4}- K_{6}) \EE_{\bxi^{[k]}}[\Delta_{\bx}^{k}] + \frac{\sigma^{2}}{b_{k}L_{F}}.
    \end{align}
    Then substituting the values of coefficients in \eqref{Eq: the coefficients of the Lyapunov function in SGDM}
    and parameter conditions \eqref{Eq: the parameters conditions of dualhash-sgdm in suppl}
    into \eqref{the descent property of Lyapunov func in SGDM first}, we derive
    \begin{align*}
             & \EE_{\bxi^{[k+1]}}\left[\Psi_m(\mathcal{Q}^{k+1})\right] - \EE_{\bxi^{[k]}}\left[\Psi_m(\mathcal{Q}^{k})\right] \nonumber                                                                                                                                                        \\
        \leq & \, -C_{\bB}\left( \EE_{\bxi^{[k+1]}}\left[\Delta_{\bB}^{k+1}\right] + \EE_{\bxi^{[k]}}\left[\Delta_{\bB}^{k}\right] \right) -C_{\bx}\left( \EE_{\bxi^{[k+1]}}\left[\Delta_{\bm{x}}^{k+1}\right]+ \EE_{\bxi^{[k]}}\left[\Delta_{\bm{x}}^{k}\right] \right) + \frac{\sigma^{2}}{b_{k}L_{F}}.
    \end{align*}Taking the full expectation of both sides leads to \eqref{Eq: the descent coefficients of potential function in SGDM},
    thereby completing the proof. 
\end{proof}
The technical proof that the coefficients of $\Psi_m$ in \eqref{Eq: the coefficients of the Lyapunov function in SGDM} and the descent constants $C_{\bB}$, $C_{\bx}$ are indeed positive is given as follows.
\begin{lemma}[Positivity of Lyapunov related coefficients]
    \label{lem:lyapunov-coef-positivity} If parameter conditions \eqref{Eq: the parameters conditions of dualhash-sgdm in suppl}
    hold, then the following inequalities are satisfied:
    \begin{align}
        % \label{Eq: the positive coefficients}
        \begin{cases}K_{5}- K_{3}- C_{3}- C_{4}\geq \nu C_{4}, \\ C_{4}- K_{6}\geq \nu C_{4},\\ C_{4}> 0,\\ C_{\bB}> 0 ,\\ C_{2},\, C_{1}> 0, \\ C_{\bx}> 0.\end{cases} \notag
    \end{align}
\end{lemma}
\begin{proof}
    For $K_{5}- K_{3}- C_{3}- C_{4}$ and $C_{4}- K_{6}$ in \eqref{the descent property of Lyapunov func in SGDM first},
    we obtain by straightforward computations by choosing some $\nu > 0$,
    {\begin{align}
        K_{5}- K_{3}- C_{3}- C_{4}
        &= \frac{1 - \alpha_{k}}{\eta_{k}}- \bar{L}_{F}- (1 + \nu)C_{4}+ \nu C_{4}\notag\\
        &= (1 - \alpha_{k})\left(\frac{1}{\eta_{k}}- \frac{\bar{L}_{F}+ (1 + \nu)C_{4}}{1 - \alpha_{k}}\right) + \nu C_{4}\notag \\
        &\geq \nu C_{4}. \notag % \label{Eq: the scaling of the delta1}
    \end{align}}
    where the last inequality follows from $\alpha_{k}= \alpha$ and
    $\frac{1}{\eta_{k}}= \frac{\bar{L}_{F}+ (1 + \nu)C_{4}}{1 - \alpha}$. Then,
    one has
    \begin{align}
        C_{4}- K_{6} & = C_{4}- \frac{\alpha}{\eta_{k}}- 2L_{F}\beta_{k}^{2}\notag                                                                                                                                   \\
                     & = \frac{1}{1 - \alpha}\left( (1 - \alpha)C_{4}- \alpha\bar{L}_{F}- (1 + \nu)C_{4}\alpha - 2L_{F}\beta^{2}(1 - \alpha) \right)\notag                                                           \\
                     & \geq \frac{1}{1 - \alpha}\left( (1 - 2\alpha - \nu)C_{4}- \alpha \left(\bar{L}_{F}+ \frac{2L_{F}\beta^{2}}{\alpha}\right) \right) + \nu C_{4}+ \frac{2L_{F}\beta^{2}\alpha}{1 - \alpha}\notag \\
                     & \geq \nu C_{4}. \notag % \label{Eq: the scaling of the delta2}
    \end{align}
    where the last inequality follows from $\alpha \in (0, \frac{1 - \nu}{2})$
    and
    $C_{4}= \frac{\bar{L}_{F}+ \frac{2L_F \beta^2}{\alpha}}{1 - 2\alpha - \nu}> 0$.
    We rearrange $\frac{1}{\eta_{k}}$ using the definition of $C_{4}$ and obtain
    $\eta_{k}= \frac{\eta}{L_{F}}$ where $\eta = \frac{(1 - 2\alpha - \nu)/2}{\frac{\bar{L}_{F}}{L_{F}}
    + \frac{(1 + \nu)\beta^{2}}{(1 - \alpha)\alpha}}$.

    Next, we verify that the coefficients of $\Psi_m$ satisfying \eqref{Eq: the coefficients of the Lyapunov function in SGDM}
    and $\tau$ satisfying \eqref{Eq: the parameters conditions of dualhash-sgdm in suppl}
    ensure that $C_{\bB},\, C_{\bx}> 0$. Let us denote $t = \tau L_{F}$. Choosing
    $\delta>0$, we have
    \begin{align*}
        C_{\bB}= \frac{L_{F}}{2}\left(\frac{2 - 6\delta}{t}- \left(3\delta + \frac{1}{\delta}\right)t- (6\delta+1) \right) \geq 0.
    \end{align*}
    If we choose $\delta = \frac{1}{6}> 0$, then
    \[
        \frac{2}{L_{F}}C_{\bB}= \frac{2}{t}- \frac{13}{2}t - 2 \geq 0.
    \]
    By selecting
    $\tau \in \left(0, 6(\frac{\sqrt{30}-2}{13})\delta \frac{1}{L_{F}}\right]$, it
    is sufficient to ensures $C_{\bB}> 0$, and consequently, $C_{2},\,C_{1}> 0$.
    Therefore, for some $\delta >0$, setting
    $\tilde{\delta}= c_{\delta}\delta^{2}$ and ensuring
    $\tau^{2}L_{F}^{2}\leq \tilde{\delta}$ guarantee positive coefficients for
    $\Psi_m$ in \eqref{Eq: the coefficients of the Lyapunov function in SGDM}
    and $C_{\bB}, \, C_{\bx}>0$. 
\end{proof}
We then derive an upper error bound on $\text{dist}^{2}(\bm{0}, \partial \calL(\bx^{k}, \bB^{k}, \bm{\Lambda}^{k}))$ based on
\begin{align}
    \label{Eq: the stationarity error bound in proof}\text{dist}^{2}(\bm{0}, \partial \calL(\bx^{k}, \bB^{k}, \bm{\Lambda}^{k})) \leq \| \bd^{k}\|^{2}=\|\bd_{\bx}^{k}\|^{2}+ \| \bd_{\bB}^{k}\|^{2}+ \| \bd_{\bm{\Lambda}}^{k}\|^{2}.
\end{align}

\begin{lemma}[Stationarity error bound]
    \label{Lemma: Stationarity error bound in SGDM} Under \Cref{Ass: main assumptions},
    for \(\forall\) $k \geq 1$, it holds that
    \begin{align}\label{Eq: the stationarity error bound in SGDM Lemma}
        \EE [\text{dist}^{2}(\bm{0}, \partial \calL(\bx^{k}, \bB^{k}, \bm{\Lambda}^{k}))] \leq S_{1}\EE[ \Delta_{\bB}^{k+1}] + S_{2}\EE[ \Delta_{\bB}^{k}] + S_{3}^{k}\EE[ \Delta_{\bm{x}}^{k+1}] + S_{4}^{k}\EE [\Delta_{\bm{x}}^{k}] + \frac{3\sigma^{2}}{b_{k}L_{F}},
    \end{align}
    where the coefficients are given by
    \begin{align}
        \label{Eq: the coefficients of stationarity error bound in SGDM}\begin{cases}S_{1}&= 4\tau^{-2}+ 12,\\ S_{2}&= 4(1 + 3 (1 + \tau L_{F})^{2}),\\ S_{3}^{k}&= 4(\eta_{k}^{-2}+ L_{F}^{2}+ 3 (\tau L_{F})^{2}),\\ S_{4}^{k}&= 6 ( L_{F}\beta_{k}^{2}+ 6 \eta_{k}^{-2}\alpha_{k}^{2}).\end{cases}
    \end{align}
\end{lemma}
% \begin{proof}
% For estimating \(\bd_{\bx}^k\), using the optimality condition of \(\bx^k\), we obtain
% \begin{align}
%     &\|\bd_{\bx}^k\|^2 \notag \\
%     =&\,\|\nabla_{\bx} F(\bx^k, \bm{B}^k) \|^2  \notag\\
%     =&\,\eta_k^{-2} \left\| \eta_k \nabla_{\bx} F(\bx^k, \bm{B}^k) + 0\right\|^2 \notag \\
%     =&\, \eta_k^{-2}  \left\|\eta_k\left( \nabla_{\bx} F(\bx^k, \bm{B}^k) - \bG^k \right) + \bm{y}^k - \bx^{k+1} \right\|^2 \notag \\
%     =&\,\eta_k^{-2} \bigg\|\eta_k\Big(
%     \nabla_{\bx} F(\bx^k, \bm{B}^k) - \nabla_{\bx} F(\bm{z}^k, \bm{B}^k) \notag \\
%     &\qquad + \nabla_{\bx} F(\bm{z}^k, \bm{B}^k)- \bG^k \Big)
%     + \bm{y}^k - \bx^k + \bx^k - \bx^{k+1} \bigg\|^2
%     \label{eq:ineq1} \\
%     \leq&\, 3 \|\nabla_{\bx} F(\bx^k, \bm{B}^k) - \nabla_{\bx} F(\bm{z}^k, \bm{B}^k)\|^2
%     + 3\eta_k^{-2}\|\bx^{k+1} - \bx^k + \bx^k - \bm{y}^k \|^2 \notag \\
%     \quad& + 3 \|\be^k\|^2
%     \label{eq:ineq2} \\
%     \leq&\, 3L_F^2 \|\bx^k - \bm{z}^k \|^2
%     + 6\eta_k^{-2} \left(\|\bx^{k+1} - \bx^k\|^2 + \|\bx^k - \bm{y}^k \|^2\right) \notag \\
%     \quad& + 3 \|\be^k\|^2,
% \end{align} where \eqref{eq:ineq1} is deduced by using the fact that \((a + b+ c)^2 \leq 3(a^2 + b^2 +c^2)\) and \eqref{eq:ineq2} is deduced by the smoothness of \( F(\bx, \bm{B})\) and \((a + b)^2 \leq 2(a^2 + b^2)\).
% Then using \cref{Pp: relations of xyz} concludes
%     \begin{align}
%     \label{Eq: dx}
%     \| \bd_{\bx}^k\| \leq  6( L_F \beta_k^2 + 2 \eta_k^{-2} \alpha_k^2) \Delta_{\bx}^{k} + 12 \eta_k^{-2} \Delta_{\bx}^{k+1} + 3 \| \be^k\|^2
%     \end{align}
%     For estimating  \(\bd^k_{\bB}\), using the optimality condition of \(\bB^k\) and the smoothness of \(F\), we obtain
% \begin{align}
%     \|\bd_{\bB}^k\|^2&= \| \nabla_{\bB} F(\bx^{k+1}, \bB^k)+  \bm{\Lambda}^k + \nabla_{\bB} F(\bx^k, \bB^k) - \nabla_{\bB} F(\bx^{k+1}, \bB^k) \|^2 \notag \\
%                     &= 4 \tau^{-2} \Delta^{k+1}_{\bB}  + 4L_F^2\Delta_{\bx}^{k+1} \label{Eq: dB},
% \end{align} where the last inequality follows from the fact \((a + b)^2 \leq 2(a^2 + b^2)\).

% For estimating \( \|\bd^k_{\bm{\Lambda}}\|^2\),  using the optimality condition of \( \bm{\Lambda}^k \) in \eqref{Eq: the optimality condition of lbda},  we obtain
% \begin{align*}
%         \|\bd^k_{\bm{\Lambda}}\|^2 &\leq 2\tau^{2} \|\bm{\Lambda}^k - \bm{\Lambda}^{k-1} \|^2+ 2\|\bB^{k} - \bB^{k-1}\|^2.
%     \end{align*} Then, substituting \eqref{Eq: dual boundedness} in \cref{Lemma: dual boundedness} into the above inequality yield that
% \begin{align}
% \label{Eq: dLambda}
% \|\bd^k_{\bm{\Lambda}} \|^2 \leq 12 \Delta_{\bB}^{k+1} + 4(1 + 3 (1 + \tau L_F )^2)\Delta_{\bB}^k + 12(\tau L_F)^2\Delta_{\bx}^{k+1} .
% \end{align}
% Then, substituting \eqref{Eq: dx}, \eqref{Eq: dB} and \eqref{Eq: dLambda} into \eqref{Eq: the stationarity error bound in proof} yields \eqref{Eq: the stationarity error bound in SGDM Lemma} and completes the proof. 
% \end{proof}
\begin{proof}
    For estimating $\| \bd_{\bx}^{k}\|^{2}$, using the optimality condition of
    $\bx^{k}$, we obtain
    \begin{align}
             & \|\bd_{\bx}^{k}\|^{2}\notag                                                                                                                                                                                                                   \\
        =    & \,\|\nabla_{\bx}F(\bx^{k}, \bm{B}^{k}) \|^{2}\notag                                                                                                                                                                                           \\
        =    & \,\eta_{k}^{-2}\left\| \eta_{k}\nabla_{\bx}F(\bx^{k}, \bm{B}^{k}) + 0\right\|^{2}\notag                                                                                                                                                       \\
        =    & \, \eta_{k}^{-2}\left\|\eta_{k}\left( \nabla_{\bx}F(\bx^{k}, \bm{B}^{k}) - \bG^{k}\right) + \bm{y}^{k}- \bx^{k+1}\right\|^{2}\notag                                                                                                           \\
        =    & \,\eta_{k}^{-2}\bigg\|\eta_{k}\Big( \nabla_{\bx}F(\bx^{k}, \bm{B}^{k}) - \nabla_{\bx}F(\bm{z}^{k}, \bm{B}^{k}) + \nabla_{\bx}F(\bm{z}^{k}, \bm{B}^{k})- \bG^{k}\Big) + \bm{y}^{k}- \bx^{k}+ \bx^{k}- \bx^{k+1}\bigg\|^{2}\label{eq:d_x_ineq1} \\
        \leq & \, 3 \|\nabla_{\bx}F(\bx^{k}, \bm{B}^{k}) - \nabla_{\bx}F(\bm{z}^{k}, \bm{B}^{k})\|^{2}+ 3\eta_{k}^{-2}\|\bx^{k+1}- \bx^{k}+ \bx^{k}- \bm{y}^{k}\|^{2}+ 3 \|\be^{k}\|^{2}\label{eq:d_x_ineq2}                                                 \\
        \leq & \, 3L_{F}^{2}\|\bx^{k}- \bm{z}^{k}\|^{2}+ 6\eta_{k}^{-2}\left(\|\bx^{k+1}- \bx^{k}\|^{2}+ \|\bx^{k}- \bm{y}^{k}\|^{2}\right) + 3 \|\be^{k}\|^{2},\label{eq:d_x_ineq3}                                                                      
    \end{align}
    where \eqref{eq:d_x_ineq1} is deduced by using the fact that
    $(a + b+ c)^{2}\leq 3(a^{2}+ b^{2}+c^{2})$ and \eqref{eq:d_x_ineq2} is deduced
    by the smoothness of $F(\bx, \bm{B})$ and $(a + b)^{2}\leq 2(a^{2}+ b^{2})$.
    Then applying \Cref{Pp: relations of xyz} to \eqref{eq:d_x_ineq3} concludes
    \begin{align}
        \label{Eq: dx}\| \bd_{\bx}^{k}\|^{2}\leq 6( L_{F}\beta_{k}^{2}+ 2 \eta_{k}^{-2}\alpha_{k}^{2}) \Delta_{\bx}^{k}+ 12 \eta_{k}^{-2}\Delta_{\bx}^{k+1}+ 3 \| \be^{k}\|^{2}.
    \end{align}
    For estimating $\|\bd^{k}_{\bB}\|^{2}$, using the $\bB$-update
    \eqref{Eq: the update of B} and the smoothness of
    $F$, we obtain
    \begin{align}
        \|\bd_{\bB}^{k}\|^{2} & = \| \nabla_{\bB}F(\bx^{k+1}, \bB^{k})+ \bm{\Lambda}^{k}+ \nabla_{\bB}F(\bx^{k}, \bB^{k}) - \nabla_{\bB}F(\bx^{k+1}, \bB^{k}) \|^{2}\notag \\
                              & \leq4 \tau^{-2}\Delta^{k+1}_{\bB}+ 4L_{F}^{2}\Delta_{\bx}^{k+1}\label{Eq: dB},
    \end{align}where the last inequality follows from the fact
    $(a + b)^{2}\leq 2(a^{2}+ b^{2})$.

    For estimating $\|\bd^{k}_{\bm{\Lambda}}\|^{2}$, using the
    $\bm{\Lambda}$-update \eqref{Eq: the update of lambda}, we obtain
\begin{align}
    \|\bd^{k}_{\bm{\Lambda}}\|^{2} & \leq 2\tau^{2}\|\bm{\Lambda}^{k}- \bm{\Lambda}^{k-1}\|^{2}+ 2\|\bB^{k}- \bB^{k-1}\|^{2}. \notag
\end{align}

    Then, substituting \eqref{Eq: dual boundedness in proof} into the above inequality yields that
    \begin{align}
        \label{Eq: dLambda}\|\bd^{k}_{\bm{\Lambda}}\|^{2}\leq 12 \Delta_{\bB}^{k+1}+ 4(1 + 3 (1 + \tau L_{F})^{2})\Delta_{\bB}^{k}+ 12(\tau L_{F})^{2}\Delta_{\bx}^{k+1}.
    \end{align}

    Finally, substituting \eqref{Eq: dx}, \eqref{Eq: dB}, and \eqref{Eq: dLambda}
    into \eqref{Eq: the stationarity error bound in proof}, and taking the full expectation
    on both sides, yields \eqref{Eq: the stationarity error bound in SGDM Lemma},
    completing the proof. 
\end{proof}

% We now establish the convergence rate of \Cref{alg:our_algorithm_sgdm} as follows.
% We set $T$ as the maximum number of iterations and denote 
%     $L_{\Delta}= \frac{\max\{S_{1}, S_{2}\} + \max\{S_{3}, S_{4}\}}{\min\{C_{\bB},
%     C_{\bx}\}}$, $C_{\sigma}= \frac{L_{\Delta}}{L_{F}}+ 3$, and
%     $\Delta_{1}= \Psi_m^{1}- \Psi_m(\mathcal{Q}^{*})$
%     are constants independent of $T$, where \(\Psi_m^{1} = \Psi_m(\calQ^1)\).
% \begin{lemma}[Convergence rate of \Cref{alg:our_algorithm_sgdm}]
%     \label{lemma: convergence rate sgdm} Under the conditions of Lemma \ref{Lemma:
%     approximate sufficient descent property in SGDM}, for any $T \geq 1$, there exists
%     $R$ uniformly selected from $\{2,\ldots,T+1\}$ such that $(\bx^{R}, \bB
%     ^{R}, \bLmbd^{R})$ satisfies \eqref{equ:esp} with
%     \begin{align}
%         \label{Eq: the convergence rate of sgdm}\mathbb{E}[\mathrm{dist}^{2}(0, \partial \mathcal{L}(\bx^{R}, \bB^{R}, \bLmbd^{R}))] \leq \frac{L_{\Delta}\Delta_{1}}{T}+ C_{\sigma}\frac{\sigma^{2}}{b_{k}}.
%     \end{align}
% \end{lemma}
\subsection{Proof of \Cref{lemma: convergence rate sgdm}}
\label{Subsubsec: convergence rate SGDM}
\begin{proof}
    From \eqref{Eq: the stationarity error bound in SGDM Lemma}, we obtain \begin{align}\label{Eq: the stationarity error bound in SGDM rearrange}\mathbb{E}[\mathrm{dist}^{2}(0, \partial \mathcal{L}(\bx^{k}, \bB^{k}, \bLmbd^{k}))] \leq ( \max\{S_{1}, S_{2}\} + \max\{S_{3}, S_{4}\}) \EE [\Delta_{\bB}^{k+1}+ \Delta_{\bB}^{k}+ \Delta_{\bx}^{k+1}+ \Delta_{\bx}^{k}] + \frac{3 \sigma^{2}}{b_{k}L_{F}}.\end{align}
    Similarly, from \eqref{Eq: the descent property of Lyapunov func in SGDM}, we
    obtain
    \[
        \EE [\Psi_m(\mathcal{Q}^{k+1})] - \EE[\Psi_m(\mathcal{Q}^{k})]
        \leq -\min \{ C_{\bB}, C_{\bx}\}\EE[\Delta_{\bB}^{k+1}+ \Delta_{\bB}^{k}+
        \Delta_{\bm{x}}^{k+1}+ \Delta_{\bm{x}}^{k}] + \frac{\sigma^{2}}{b_{k}L_{F}},
    \]
    which yields that
    \[
        \EE[\Delta_{\bB}^{k+1}+ \Delta_{\bB}^{k}+ \Delta_{\bm{x}}^{k+1}+ \Delta_{\bm{x}}
        ^{k}] \leq \frac{1}{\min \{ C_{\bB}, C_{\bx}\}}\left( \EE [\Psi_m(\mathcal{Q}
        ^{k})] - \EE[\Psi_m(\mathcal{Q}^{k+1})] + \frac{\sigma^{2}}{b_{k}L_{F}}
        \right)
    \]
    Substituting this inequality into \eqref{Eq: the stationarity error bound in SGDM rearrange},
    summing over $k = 1, \cdots,T$, and dividing by $T$, we obtain
    \[
        \frac{1}{T}\sumT \mathbb{E}[\mathrm{dist}^{2}(0, \partial \mathcal{L}(\bx
        ^{k}, \bB^{k}, \bLmbd^{k}))] \leq \frac{L_{\Delta}}{T}\sumT
        \left(\EE [\Psi_m(\mathcal{Q}^{k})] - \EE [\Psi_m
        (\mathcal{Q}^{k+1})] \right) + ( 3 + L_{\Delta}) \frac{\sigma^{2}}{b_{k}},
    \], which yields the conclusion. Denoting
    $a_{k}= \mathbb{E}[\mathrm{dist}^{2}(0, \partial \mathcal{L}(\bx^{k}, \bB
    ^{k}, \bLmbd^{k}))]$, we have
    \[
        \frac{1}{T}\sumT a_{k}\leq \frac{L_{\Delta}\Delta_{1}}{T}+ ( 3 + L_{\Delta}
        ) \frac{\sigma^{2}}{b_{k}}.
    \]
    By uniformly selecting $R' \in \{1, \ldots, T\}$ and letting $R = R' + 1$,
    we obtain
    \[
        \mathbb{E}[\mathrm{dist}^{2}(0, \partial \mathcal{L}(\bx^{R}, \bB^{R}
        , \bLmbd^{R}))] \leq \frac{L_{\Delta}\Delta_{1}}{T}+ ( 3 + L_{\Delta}
        ) \frac{\sigma^{2}}{b_{k}},
    \] and complete the proof.
\end{proof}
% From \eqref{Eq: the convergence rate of sgdm}, to achieve $\calO(T^{-1})$
% convergence rate, we need to set $b_{k}= \Theta (T)$. Next, we set
% $b_{k}= b = c_{b}T$ for all $k \geq 1$ in this way and establish the complexity
% results of \Cref{alg:our_algorithm_sgdm} with the parameter conditions \eqref{Eq: the parameters conditions of dualhash-sgdm in suppl}.
\subsection{Proof of \Cref{Theorem: the oracle complexity in SGDM}}
\label{Subsec: proof theorem SGDM}
% \setcounter{proof}{0}
\begin{proof}
    When $b_{k}= b = c_{b}T$, we have
    \[
        \mathbb{E}[\mathrm{dist}^{2}(0, \partial \mathcal{L}(\bx^{R}, \bB^{R}
        , \bLmbd^{R}))] \leq \frac{1}{T}\left( L_{\Delta}\Delta_{1}+ C_{\sigma}
        \frac{\sigma^{2}}{c_{b}}\right).
    \]
    With the conditions on $T$ in \eqref{Eq: the iterate number upper bound of algorithm SGDM},
    this gives
    $\mathbb{E}[\mathrm{dist}^{2}(0, \partial \mathcal{L}(\bx^{R}, \bB^{R}
    , \bLmbd^{R}))] \leq \varepsilon^{2}$. Then, the oracle complexity
    of SFO is
    \[
        \text{Complexity}= (T-1)b_{k}=\calO(T^{2}) = \calO(\varepsilon^{-4}).
    \]
\end{proof}

\section{Complexity analysis of DualHash-StoRM (\texorpdfstring{\Cref{alg:our_algorithm_storm}}{Algorithm 2})}
\label{Appendix: the complexity analysis of DualHash-StoRM}
We adopt the same notation as in \Cref{Appendix: the complexity analysis of DualHash-StoM}. 
\subsection{Proof of \Cref{Lemma: approximate sufficient descent property storm}}
\label{Subsubsec: proof descent STORM}
\begin{proof}
    Substituting the one-iteration progress \eqref{Eq: the recursive relation of L on x}
    and \eqref{Eq: the recursive relation of L on B and Lambda} into \eqref{Eq: One-iteration progress decomposition},
    we obtain
    \begin{align}
        % \label{Eq: the recursive relation of L under VR}
        & \calL (\bx^{k+1}, \bB^{k+1}, \bLmbd^{k+1}) - \calL (\bx^{k}, \bB^{k}, \bLmbd^{k}) \notag                                                                                                        \\
        \leq \,                                          & \frac{s}{2}\| \be^{k}\|^{2}- (-\frac{1}{s}+ 2\eta_{k}^{-1}- L_{F}- K_{3}) \Delta_{\bx}^{k+1}+ K_{1}\Delta_{\bB}^{k+2}+ K_{4}\Delta_{\bx}^{k+2}- K_{2}\Delta_{\bB}^{k+1}+ K_{3}\Delta_{\bB}^{k}. \notag
    \end{align}
    where setting $s = \eta_{k}$ yields
    \begin{align}
             & \calL (\bx^{k+1}, \bB^{k+1}, \bLmbd^{k+1}) - \calL (\bx^{k}, \bB^{k}, \bLmbd^{k}) \nonumber                                                                                                                                          \\
        \leq & \,\frac{\eta_{k}}{2}\| \be^{k}\|^{2}- (\eta_{k}^{-1}- L_{F}- K_{3}) \Delta_{\bx}^{k+1}+ K_{1}\Delta_{\bB}^{k+2}+ K_{4}\Delta_{\bx}^{k+2}- K_{2}\Delta_{\bB}^{k+1}+ K_{3}\Delta_{\bB}^{k}.\label{Eq: One-iteration progress in storm}
    \end{align}
    Applying the definition of $\Psi_r$ from \eqref{Eq: the Lyapunov function in STORM}
    into \eqref{Eq: One-iteration progress in storm}, we rearrange the above
    inequality
    \begin{align}
        \label{the descent property of Lyapunov func in STORM first} & \Psi_r(\mathcal{Q}^{k+1}) - \Psi_r(\mathcal{Q}^{k}) \notag                                                                                         \\
        \leq                                                         & - (C_{1}- K_{1}) \Delta_{\bm{B}}^{k+2}-(K_{2}- C_{1}- C_{2}) \Delta_{\bm{B}}^{k+1}- (C_{2}- K_{3}) \Delta_{\bm{B}}^{k}- (C_{3}-K_{4}) \Delta_{\bx}^{k+2}\notag \\
                                                                     & -(\eta_{k}^{-1}- L_{F}- K_{3}- C_{3}) \Delta_{\bx}^{k+1}- (C_{4}- K_{6}) \Delta_{\bx}^{k}+ \frac{\eta_{k}}{2}\| \be^{k}\|^{2}.
    \end{align}
    By substituting the Lyapunov function coefficients \eqref{Eq: the coefficients of the Lyapunov function in SGDM}
    into \eqref{the descent property of Lyapunov func in STORM first}, we obtain
    \eqref{Eq: the descent property of Lyapunov func in storm}. With $\tau$ satisfying
    \eqref{Eq: the parameters conditions of dualhash-sgdm in suppl}, \Cref{lem:lyapunov-coef-positivity}
    guarantees that $C_{\bB}> 0$. 
\end{proof}

\subsection{Proof of \Cref{Lemma: Weighted error accumulation}}
\label{Subsubsec: proof Weighted error accumulation}
\begin{proof}
    Under \eqref{Eq: the parameters assumptions in storm appendix}, we have a lower
    bound for $C_{\bx}^{k}$ with
    \begin{align}
        \label{Eq: the bound coefficient of x}C_{\bx}^{k}\geq \frac{1}{2\eta_{1}}>0 .
    \end{align}
    Substituting this inequality into the second inequality of \eqref{Eq: bound B and x in storm}
    and taking the conditional expectation on both sides yields
    \begin{align}
        \label{Eq: scaling recursive}\EE_{\bxi^{[k+1]}}\left[\Psi_r(\calQ^{k+1})\right] - \EE_{\bxi^{[k]}}\left[\Psi_r(\calQ^{k})\right] \leq - \frac{1}{2\eta_{1}}\EE_{\bxi^{[k+1]}}\left[ \Delta_{\bx}^{k+1}\right] + \frac{\eta_{k}}{2}\EE_{\bxi^{[k]}}\left[ \|\be^{k}\|^{2}\right].
    \end{align}
    From \Cref{Lemma: Recursive bound on gradient error}, we obtain
    \begin{align}
        \EE_{\bxi^{[k+1]}}\left[\| \be^{k+1}\|^{2}\right] \leq \left(1 - \rho_{k}\right)\EE_{\bxi^{[k]}}\left[\| \be^{k}\|^{2}\right] + \frac{2 \rho_{k}^{2}\sigma^{2}}{|\mathcal{J}_{k+1}|}+ 4(1-\rho_{k})L_{F}^{2} \EE_{\bxi^{[k+1]}}\left[ \Delta_{\bx}^{k+1} \right], \nonumber
    \end{align}
    where the inequality follows from $0 < \rho_{k}<1$ and  $| \calJ_{k+1}|
    \geq 1$, and this implies
    \begin{align}
        &- \EE_{\bxi^{[k+1]}}\left[ \Delta_{\bx}^{k+1} \right] \notag \\
        \leq&\, \frac{1}{4(1-\rho_{k})^{2}L_{F}^{2}}\left(\EE_{\bxi^{[k]}}\left[\| \be^{k}\|^{2}\right] - \EE_{\bxi^{[k+1]}}\left[\| \be^{k+1}\|^{2}\right] -\rho_k \EE_{\bxi^{[k]}}\left[\| \be^{k}\|^{2}\right]+ \frac{2 \rho_{k}^{2}\sigma^{2}}{|\mathcal{J}_{k+1}|}\right). \nonumber
    \end{align}
    Substituting this inequality into \eqref{Eq: scaling recursive}, we obtain
    \begin{align}
             & \EE_{\bxi^{[k+1]}}\left[\Psi_r(\calQ^{k+1})\right] - \EE_{\bxi^{[k]}}\left[\Psi_r(\calQ^{k})\right] \notag     \\
             \leq & \frac{1}{8(1-\rho_k)\eta_1L^2_F}\left(\EE_{\bxi^{[k]}}\left[\| \be^{k}\|^{2}\right] - \EE_{\bxi^{[k+1]}}\left[\| \be^{k+1}\|^{2}\right] -\rho_k \EE_{\bxi^{[k]}}\left[\| \be^{k}\|^{2}\right]+ \frac{2 \rho_{k}^{2}\sigma^{2}}{|\mathcal{J}_{k+1}|}\right) + \frac{\eta_{k}}{2}\EE_{\bxi^{[k]}}\left[ \|\be^{k}\|^{2}\right]  \notag \\
        \leq & \frac{1}{8\eta_{1}L_{F}^{2}}\left(\EE_{\bxi^{[k]}}\left[\| \be^{k}\|^{2}\right] - \EE_{\bxi^{[k+1]}}\left[\| \be^{k+1}\|^{2}\right] + \frac{2 \rho_{k}^{2}\sigma^{2}}{|\mathcal{J}_{k+1}|}\right) - \left( \frac{\rho_{k}}{8\eta_{1}L_{F}^{2}}- \frac{\eta_{k}}{2}\right)\EE_{\bxi^{[k]}}\left[\| \be^{k}\|^{2}\right] \notag \\
        &+ \frac{\rho_{k}^{2}\sigma^{2}}{4\eta_{1}L_{F}^{2}| \calJ_{k+1}|}, \label{Eq: scaling recursive 2}
    \end{align}where the second inequality follows from $ 1 - \rho_{k}\leq 1$.
    From \eqref{Eq: the parameters assumptions in storm appendix}, it follows that
    $\frac{\rho_{k}}{8\eta_{1}L_{F}^{2}}- \frac{\eta_{k}}{2}\geq \frac{\rho_{k}}{16\eta_{1}L_{F}^{2}}$.
    Applying this inequality to \eqref{Eq: scaling recursive 2} yields the following
    expected descent property of the Lyapunov function $\Psi_r$
    \begin{align}
        \label{Eq: One-iteration Expected descent} & \EE_{\bxi^{[k+1]}}\left[ \Psi_r(\calQ^{k+1}) + \frac{1}{8\eta_{1}L_{F}^{2}}\| \be^{k+1}\|^{2}\right] - \EE_{\bxi^{[k]}}\left[ \Psi_r(\calQ^{k}) + \frac{1}{8\eta_{1}L_{F}^{2}}\| \be^{k}\|^{2}\right] \notag \\
        \leq                                       & - \frac{\rho_{k}}{16\eta_{1}L_{F}^{2}}\EE_{\bxi^{[k]}}[\| \be^{k}\|^{2}] + \frac{\rho_{k}^{2}\sigma^{2}}{4\eta_{1}L_{F}^{2}| \calJ_{k+1}|}.
    \end{align}
    Finally, summing up \eqref{Eq: One-iteration Expected descent} from $k=1$ to
    $T$ and dividing by $T$ yields
    \[
        \frac{1}{T}\sumT \rho_{k}\EE_{\bxi^{[k]}}[\| \be^{k}\|^{2}]\leq \frac{16\eta_{1}L_{F}^{2}\tilde{\Delta}_{1}}{T}
        + \frac{2\EE_{\bxi^{[1]}}[\| \be^{1}\|^{2}]}{T}+ \frac{1}{T}\sum\limits_{k=2}
        ^{T}\frac{4 \rho_{k}^{2}\sigma^{2}}{|\calJ_{k+1}|},
    \], yielding \eqref{Eq: weighted cumulative error bound} and completes the proof. 
\end{proof}

We provide a recursive bound on the gradient error vector $\{\be^{k}\}$ with analysis comparable to those in \cite[Lemma 2]{xu2023momentum} and \cite[Lemma 8]{shi2025momentum}.
\begin{lemma}[Recursive bound on gradient error]
    \label{Lemma: Recursive bound on gradient error} Under \Cref{Ass: main assumptions},
    for $\forall \, k \geq 1$, it holds that
    \begin{align}
        % \label{Eq: recursive bound on gradient error}
        \EE_{\bxi^{[k+1]}}[\| \be^{k+1}\|^{2}] \leq \left(1 - \rho_{k}\right)^{2}\EE_{\bxi^{[k]}}[\| \be^{k}\|^{2}] + \frac{2 \rho_{k}^{2}\sigma^{2}}{|\mathcal{J}_{k+1}|}+ \frac{4(1 - \rho_{k})^{2}L_{F}^{2}}{|\mathcal{J}_{k+1}|}\EE_{\bxi^{[k+1]}}\left[\Delta_{\bx}^{k+1}\right]. \notag
    \end{align}
\end{lemma}
% \begin{remark}
%     This recursive bound reveals three key components affecting gradient error:
%     (i) historical error decay (controlled by $1-\rho_{k}$), (ii) variance from current
%     stochastic gradient estimation (proportional to $\rho_{k}^{2}$), and (iii)
%     error propagation from recent updates (scaling with $(1-\rho_{k})^{2}\Delta_{\bx}
%     ^{k+1}$). The momentum parameter $\rho_{k}$ creates a trade-off between these
%     components.
% \end{remark}
We now establish a weighted error accumulation bound
$\frac{1}{T}\sum_{k=1}^{T}\rho_{k}\EE_{\bxi^{[k]}}\left[\|\be^{k}\|^{2}\right]$
for the variance-reduced gradient estimator in \Cref{alg:our_algorithm_storm}, 
with the proof given in \Cref{Subsubsec: proof Weighted error accumulation}. 
For this analysis, we require the parameters $\eta_{k}$, $\rho_{k}$, $\tau$ 
and the batch size $b_k$ to satisfy:
\begin{align}
    \label{Eq: the parameters assumptions in storm appendix}
    \eta_{k+1}\leq \eta_{k}, \quad \eta_{k}\tilde{L}_{F}\leq \frac{1}{2}, \quad 
    0 < 8\eta_{1}\eta_{k}L_{F}^{2}\leq \rho_{k}\leq1, \quad 
    \tau \leq \frac{\sqrt{\tilde{\delta}}}{L_{F}}, \quad
    b_k = \begin{cases} b_1 & k = 1, \\ b & k \geq 2. \end{cases}
\end{align}
\begin{lemma}[Weighted error accumulation bound]
    \label{Lemma: Weighted error accumulation} Under \Cref{Ass: main assumptions}
    and parameter conditions \eqref{Eq: the parameters assumptions in storm appendix}, for $\forall$ $T \geq 1$, the following weighted cumulative error bound
    holds:
    \begin{align}
        \label{Eq: weighted cumulative error bound}\frac{1}{T}\sum_{k=1}^{T}\rho_{k}\EE_{\bxi^{[k]}}\left[ \| \be^{k}\|^{2}\right]\leq \frac{16 \eta_{1}L_{F}^{2}\tilde{\Delta}_{1}}{T}+ \frac{2\sigma^{2}}{Tb_{1}}+ \frac{1}{T}\sum_{k=2}^{T}\frac{4\rho_{k}^{2}\sigma^{2}}{b}.
    \end{align} where $\tilde{\Delta}_{1}= \Psi_r^{1}- \Psi_r^{*}$, $\Psi_r^{1}= \Psi_r(\calQ^{1})$ and $\Psi_r^{*}$ is the lower bound
of $\Psi_r$.
\end{lemma}
We derive an upper stationarity error bound for the iterates generated by \Cref{alg:our_algorithm_storm}.
\begin{lemma}[Stationarity error bound]
    \label{Lemma: Stationarity error bound} Under \Cref{Ass: main assumptions} (ii),
    for $\forall$ $k \geq 1$, it holds that
    \begin{align}
        \text{dist}^{2}(\bm{0}, \partial \calL(\bx^{k}, \bB^{k}, \bm{\Lambda}^{k})) \leq S_{1}\Delta_{\bB}^{k+1}+ S_{2}\Delta_{\bB}^{k}+ S_{3}^{k}\Delta_{\bm{x}}^{k+1}+ 2 \| \be^{k}\|^{2}, \label{Eq: the stationarity error bound in Lemma in storm}
    \end{align}
    where the coefficients are defined as \eqref{Eq: the coefficients of stationarity error bound in SGDM}.
    % \begin{align}
    % \label{Eq: the  coefficients of stationarity error bound in storm}
    % \begin{cases}
    %     S_1 &= 4\tau^{-2} + 12,\\
    %     S_2 &=  4(1 + 3 (1 + \tau L_F )^2),\\
    %     S_3^k &= 4(\eta_k^{-2} + L_F^2+ 3 (\tau L_F)^2).
    % \end{cases}
    % \end{align}
\end{lemma}
\begin{proof}
    This proof follows a similar approach to that of \Cref{Lemma: Stationarity error bound in SGDM} 
in \Cref{Appendix: the complexity analysis of DualHash-StoRM}. The key difference lies in 
the estimation of $\| \bd_{\bx}^{k}\|^{2}$. To estimate $\| \bd_{\bx}^{k}\|^{2}$, we use 
the optimality condition of $\bx^{k+1}$ from \eqref{Eq: the update of x in STORM} to obtain
    \begin{align}
        \|\bd_{\bx}^{k}\|^{2}= & \,\eta_{k}^{-2}\left\| \eta_{k}\nabla_{\bx}F(\bx^{k}, \bm{B}^{k}) + 0\right\|^{2}\notag                                       \\
        =                      & \, \eta_{k}^{-2}\left\|\eta_{k}\left( \nabla_{\bx}F(\bx^{k}, \bm{B}^{k}) - \bD^{k}\right) + \eta_{k}\bD^{k}\right\|^{2}\notag \\
        =                      & \, \eta_{k}^{-2}\left\|\eta_{k}\be^{k}+ (\bx^{k}- \bx^{k+1})\right\|^{2}\notag                                                \\
        \leq                   & \, 4\eta_{k}^{-2}\Delta_{\bx}^{k+1}+ 2 \| \be^{k}\|^{2}.\label{Eq: dx in storm}
    \end{align}
    Then, substituting \eqref{Eq: dx in storm}, \eqref{Eq: dB}, and \eqref{Eq: dLambda}
    into \eqref{Eq: the stationarity error bound in proof} yields \eqref{Eq: the stationarity error bound in Lemma in storm}
    and completes the proof. 
\end{proof}
% \begin{align}
%     \begin{cases}
%         &\bd_{\bx}^k =  \nabla_{\bx} F(\bx^k, \bB^k),\\
%         &\bd_{\bB}^k = \nabla_{\bB} F(\bx^k, \bB^k) + \bm{\Lambda}^k, \\
%         &\bd_{\bm{\Lambda}}^k = \bB^k - \calG_{\bm{\Lambda}}^k.
%     \end{cases}
% \end{align} where \(\calG_{\bm{\Lambda}}^k \in \partial h^*(\bm{\Lambda}^k)\).
% Therefore,
% \begin{align}
% \label{Eq: the stationarity error bound in proof}
% \text{dist}^2(\bm{0}, \partial \calL(\bx^k, \bB^k, \bm{\Lambda}^k)) \leq \|\bd_{\bx}^k \|^2 + \| \bd_{\bB}^k\|^2 + \| \bd_{\bm{\Lambda}}^k\|^2.
% \end{align}
%     For estimating  \(\bd^k_{\bB}\), using the optimality condition of \(\bB^k\) and the smoothness of \(F\), we obtain
% \begin{align}
%     \|\bd_{\bB}^k\|^2&= \| \nabla_{\bB} F(\bx^{k+1}, \bB^k)+  \bm{\Lambda}^k + \nabla_{\bB} F(\bx^k, \bB^k) - \nabla_{\bB} F(\bx^{k+1}, \bB^k) \|^2 \notag \\
%                     &= 4 \tau^{-2} \Delta^{k+1}_{\bB}  + 4L_F^2\Delta_{\bx}^{k+1},
% \end{align} where the last inequality follows from the fact \((a + b)^2 \leq 2(a^2 + b^2)\).

% For estimating \( \|\bd^k_{\bm{\Lambda}}\|^2\),  using the optimality condition of \( \bm{\Lambda}^k \) in \eqref{Eq: the optimality condition of lbda},  we obtain
% \begin{align*}
%         \|\bd^k_{\bm{\Lambda}}\|^2 &\leq 2\tau^{2} \|\bm{\Lambda}^k - \bm{\Lambda}^{k-1} \|^2+ 2\|\bB^{k} - \bB^{k-1}\|^2.
%     \end{align*} Then, substituting \eqref{Eq: dual boundedness} in \cref{Lemma: dual boundedness} into the above inequality yield that
% \begin{align}
% \|\bd^k_{\bm{\Lambda}} \|^2 \leq 12 \Delta_{\bB}^{k+1} + 4(1 + 3 (1 + \tau L_F )^2)\Delta_{\bB}^k + 12(\tau L_F)^2\Delta_{\bx}^{k+1} .
% \end{align}

Selecting appropriate parameters is crucial as the ultimate oracle complexities depend on these choices.
\subsection{Proof of averaged stationarity error bound for STORM}
\label{subsubsec:proof-avg-stationarity-storm}
We now establish the convergence rate of \Cref{alg:our_algorithm_storm} by taking constant parameters, i.e., $\eta_{k}\equiv \eta_{1}$ and $\rho_{k}\equiv \rho_{1}$ for $k \geq 1$. The first result is to estimate the averaged expected stationarity error.
\begin{lemma}
    \label{lemma: the averaged stationarity error bound} 
    Under the conditions of
    \Cref{Lemma: Weighted error accumulation} , it holds that for $\forall$
    $T \geq 1$,
    \begin{align}
    \frac{1}{T}\sumT \EE [\text{dist}^{2}(\bm{0}, \partial \calL(\bx^{k}, \bB^{k}, \bm{\Lambda}^{k}))]        
        \leq & \, \frac{\tilde{\Delta}_{1}}{T}\left( \frac{\tilde{S}_{4}}{C_{\bB}}+ \frac{2(C^1_{\delta}+ 4)}{\eta_{1}}+ \frac{(6+C^1_{\delta}+ \frac{\tilde{S}_4\eta_1}{2C_{\bB}})16\eta_{1}L_{F}^{2}}{\rho_{1}}\right) \notag \\
        &+ \left( 6 + C^1_{\delta}+ \frac{\tilde{S}_{4}\eta_{1}}{2C_{\bB}}\right)\left(\frac{2\sigma^{2}}{b_{1}\rho_{1}T}+ \frac{4\rho_{1}\sigma^{2}(T-1)}{bT}\right), \label{Eq: the averaged stationarity error bound in lemma}
    \end{align}
\end{lemma} where $\tilde{S}_4= \max\{S_{1}, S_{2}\}$ and $C^1_{\delta}= 1 + \frac{3\tilde{\delta}}{\tilde{L}_{F}^{2}}$.
\begin{proof}
    Applying $\eta_{k}\leq \eta_{1}$ from \eqref{Eq: the parameters assumptions in storm appendix}
    to \eqref{Eq: the stationarity error bound in Lemma in storm} in \Cref{Lemma:
    Stationarity error bound} yields
    \begin{align}
        \label{Eq: the stationarity error bound in proof storm}\text{dist}^{2}(\bm{0}, \partial \calL(\bx^{k}, \bB^{k}, \bm{\Lambda}^{k})) \leq & \,\frac{S_{4}}{C_{\bB}}\left( \Psi_r(\calQ^{k}) - \Psi_r(\calQ^{k+1}) \right) + S_{3}^{k}\Delta_{\bm{x}}^{k+1}+(2 + \frac{S_{4}\eta_{1}}{2C_{\bB}})\| \be^{k}\|^{2}.
    \end{align}
    For the second term in \eqref{Eq: the stationarity error bound in proof storm},
    by recalling the second inequality in \eqref{Eq: bound B and x in storm}, we
    obtain
    \[
        \Delta_{\bm{x}}^{k+1}\leq \frac{1}{ C_{\bx}^{k}}\left( \Psi_r(\mathcal{Q}
        ^{k}) - \Psi_r(\mathcal{Q}^{k+1}) +\frac{\eta_{k}}{2}\| \be^{k}\|^{2}
        \right),
    \] where we use the positive value of $C_{\bx}^{k}$ from \eqref{Eq: the bound coefficient of x}. Therefore, substituting the inequality above and the definition of
    $S_{3}^{k}$ from \eqref{Eq: the coefficients of stationarity error bound in SGDM},
    into $S_{3}^{k}\Delta_{\bx}^{k+1}$, one has
\begin{align}
\label{Eq: the scaling for x sequence}
&S_{3}^{k}\Delta_{\bx}^{k+1} \notag \\
\leq &\,\frac{S_{3}^{k}}{C_{\bx}^{k}}\left( \Psi_r(\mathcal{Q}^{k}) - \Psi_r(\mathcal{Q}^{k+1}) +\frac{\eta_{k}}{2}\| \be^{k}\|^{2}\right)\notag    \\
\leq &\, \frac{4 [ 1 + (1 + 3\tau^{2})(\eta_{k}L_{F})^{2}]}{\eta_{k}(1 - \eta_{k}\tilde{L}_{F})}\left( \Psi_r(\mathcal{Q}^{k}) - \Psi_r(\mathcal{Q}^{k+1}) +\frac{\eta_{k}}{2}\| \be^{k}\|^{2}\right) \notag                                                                                                                                    \\
\leq &\, \frac{4 [ 1 + (1 + \frac{3\tilde{\delta}}{L^2_F})(\eta_{k}L_{F})^{2}]}{\eta_{k}(1 - \eta_{k}\tilde{L}_{F})}\left( \Psi_r(\mathcal{Q}^{k}) - \Psi_r(\mathcal{Q}^{k+1}) + \frac{\eta_{k}}{2}\| \be^{k}\|^{2}\right) \notag                                                                                                               \\
\leq &\, \frac{4 [ 1 + (1 + \frac{3\tilde{\delta}}{\tilde{L}^2_F})(\eta_{k}\tilde{L}_{F})^{2}]}{\eta_{1}(1 - \eta_{k}\tilde{L}_{F})}\left( \Psi_r(\mathcal{Q}^{k}) - \Psi_r(\mathcal{Q}^{k+1})\right) + \frac{2 [ 1 + (1 + \frac{3\tilde{\delta}}{\tilde{L}^2_F})(\eta_{k}\tilde{L}_{F})^{2}]}{1 - \eta_{k}\tilde{L}_{F}}\| \be^{k}\|^{2}\notag \\
% &\leq \frac{4 [ 1 + C_{\delta}(\eta_k\tilde{L}_F)^2]}{\eta_k (1 - \eta_k\tilde{L}_F)} \left( \Psi_r(\mathcal{Q}^k) - \Psi_r(\mathcal{Q}^{k+1}) \right)  +  \frac{2[ 1 + C_{\delta}(\eta_k\tilde{L}_F)^2]}{\eta_1 (1 - \eta_k\tilde{L}_F)}\| \be^k\|^2\notag \\
\leq &\,\frac{4 [ 1 + C_{\delta}(\eta_{k}\tilde{L}_{F})^{2}]}{\eta_{1}(1 - \eta_{k}\tilde{L}_{F})}\left( \Psi_r(\mathcal{Q}^{k}) - \Psi_r(\mathcal{Q}^{k+1})\right) + \frac{2 [ 1 +C_{\delta}(\eta_{k}\tilde{L}_{F})^{2}]}{1 - \eta_{k}\tilde{L}_{F}}\| \be^{k}\|^{2}\notag                                                                     \\
\leq &\, \frac{2(C_{\delta}+ 4)}{\eta_{1}}\left( \Psi_r(\mathcal{Q}^{k}) - \Psi_r(\mathcal{Q}^{k+1})\right) + (C_{\delta}+ 4) \| \be^{k}\|^{2},
\end{align} where the third inequality follows from $\tau^{2}\leq \tilde{\delta}/ L_{F}^{2}$,
    the fourth inequality is derived from $\tilde{L}_{F}\geq L_{F}$ and
    $\eta_{k}\equiv \eta_{1}$, and the last inequality holds because $(1 + cx^{2}
    ) / (1 - x) \leq 2 + \frac{c}{2}$ for $0 < x \leq 1/2$ and constant $c >1$ under
    parameter conditions \eqref{Eq: the parameters assumptions in storm appendix}.

    On the other hand, by taking the conditional expectation on both sides of \eqref{Eq: the stationarity error bound in proof storm}
    and averaging over $T$, we obtain
    \begin{align}
             & \frac{1}{T}\sumT \EE_{\bxi^{[k]}}\left[\text{dist}^{2}(\bm{0}, \partial \calL(\bx^{k}, \bB^{k}, \bm{\Lambda}^{k}))\right] \notag                                                                                                                                                                      \\
        \leq & \,\frac{S_{4}}{C_{\bB}T}\sumT \left( \EE_{\bxi^{[k]}}\left[\Psi_r(\calQ^{k})\right] - \EE_{\bxi^{[k+1]}}\left[\Psi_r(\calQ^{k+1}) \right] \right)+ \frac{1}{T}\sumT S_{3}^{k}\EE_{\bxi^{[k]}}\left[ \Delta_{\bm{x}}^{k+1}\right] \notag                                                     \\
             & +\frac{1}{T}\sumT \left( 2 + \frac{S_{4}\eta_{1}}{2C_{\bB}}\right) \EE_{\bxi^{[k]}}\left[\| \be^{k}\|^{2}\right] \notag                                                                                                                                                                                 \\
        \leq & \frac{\frac{S_4}{C_B}\tilde{\Delta}_{1}}{ T}+ \frac{1}{T}\sumT S_{3}^{k}\EE_{\bxi^{[k+1]}}\left[ \Delta_{\bm{x}}^{k+1}\right]+\left( 2 + \frac{S_{4}\eta_{1}}{2C_{\bB}}\right)\frac{1}{T}\sumT \EE_{\bxi^{[k]}}\left[\| \be^{k}\|^{2}\right].\label{Eq: the averaged stationarity error bound in proof}
    \end{align}

    Then substituting \eqref{Eq: the scaling for x sequence} into \eqref{Eq: the averaged stationarity error bound in proof}
    yields
    \begin{align}
             & \frac{1}{T}\sumT \EE_{\bxi^{[k]}}[\text{dist}^{2}(\bm{0}, \partial \calL(\bx^{k}, \bB^{k}, \bm{\Lambda}^{k}))] \notag                                                                                                                                                     \\
        \leq & \frac{ \left( \frac{S_4}{C_B}+ \frac{2(C_{\delta}+ 4)}{S_4\eta_1}\right)\tilde{\Delta}_{1}}{ T}+\left( 6+ \frac{S_{4}\eta_{1}}{2C_{\bB}}+ C_{\delta}\right)\frac{1}{T}\sumT \EE_{\bxi^{[k]}}[\| \be^{k}\|^{2}].\label{Eq: the averaged stationarity error bound in proof 2}
    \end{align}

    Since $\rho_{k}\equiv \rho_{1}$ for $k \geq 1$, then it follows from \eqref{Eq: weighted cumulative error bound}
    that
    \begin{align}
        \label{Eq: the averaged cumulative error bound}\frac{1}{T}\sumT \EE_{\bxi^{[k]}}[\| \be^{k}\|^{2}]\leq \frac{16 \eta_{1}L_{F}^{2}\tilde{\Delta}_{1}}{\rho_{1}T}+ \frac{2\sigma^{2}}{b_{1}T}+ \frac{4\rho_{1}\sigma^{2}(T-1)}{bT}
    \end{align}
    Then, substituting \eqref{Eq: the averaged cumulative error bound} into \eqref{Eq: the averaged stationarity error bound in proof 2}
    and taking the full expectation on both sides yields \eqref{Eq: the averaged stationarity error bound in lemma}.
\end{proof}

\subsection{Proof of \Cref{Lemma: convergence rate of storm}}
\label{Subsubsec: convergence rate STORM}
\begin{proof}
    Substituting the parameter conditions \eqref{Eq: the parameters conditions in dualhash-storm} into \eqref{Eq: the averaged stationarity error bound in lemma} from \Cref{lemma: the averaged stationarity error bound}, and applying the bound $\rho_k \equiv \frac{8\rho\eta^2}{T^{2/3}}$, $b_1 = c_b T^{1/3}$, $b_k = b$ for $k \geq 2$, we obtain \eqref{Eq: the convergence rate of storm} by selecting $R$ uniformly from $\{2,\ldots,T+1\}$.
\end{proof}
With this parameter choice, we proceed to analyze the oracle complexity of \Cref{alg:our_algorithm_storm}.
\subsection{Proof of \Cref{Theorem: the oracle complexity in STORM}}
\label{Subsubsec: proof theorem STORM}
\begin{proof}
    When $b_{1}= c_{b}T^{1/3}$, we have
    \[
        \EE[\text{dist}^{2}(0, \partial \calL(\bx^{R}, \bB^{R}, \bm{\Lambda}^{R})
        ] \leq \frac{1}{T^{2/3}}\left(\tilde{L}_{\Delta}\tilde{\Delta}_{1}+ \tilde{C}_{\sigma}
        \left( \frac{\sigma^{2}}{ 4\rho \eta^{2}c_{b}}+ \frac{32\rho\eta^{2}\sigma^{2}}{b}
        \right) \right).
    \]
    With the conditions on $T$ in \eqref{Eq: the iterate number upper bound of algorithm VR},
    this gives
    $\EE[\text{dist}^{2}(0, \partial \calL(\bx^{R}, \bB^{R}, \bm{\Lambda}^{R}))]
    \leq \varepsilon^{2}$. Then, the oracle complexity is
    \[
        \text{Complexity}= b_{1}+ (T-1)b = \calO(T^{\frac{1}{3}}) + \calO(T) = \calO
        (T) = \calO(\varepsilon^{-3}),
    \]
    which completes the proof. 
\end{proof}
% \[
% \eta_k = \frac{\eta}{L_F} \frac{1}{T^l}, \quad (L_F + K_3) \eta_k \leq \frac{1}{2}
% \] where \(\rho\) and \(\eta\) are independent of \(T\).
% Let us \(l = \frac{1}{3}\) and \(q = \frac{2}{3}\),
% we want \(\frac{1}{T}\sum_{k=1}^T\text{dist}^2(\bm{0}, \partial \calL^k) = \mathcal{O}(T^{-\frac{2}{3}})\)
% We denoteare constant independent of \(T\) and \(\Delta_1 = \Psi_r(\calQ^1) - \Psi_r^* \) and the definitions of the \(\Psi_r\) and \(\calQ\) are detailed in Appendix \ref{Appendix: the complexity analysis of DualHash-StoRM}.

% Let us discuss how to select appropriate parameters.  To ensure the parameter setting \eqref{Eq: the parameters conditions in dualhash-storm}, we set the parameters as follows
% \[
% \eta_k = \frac{\eta}{T^l\tilde{L}_{F}},\quad \rho_k = \frac{8\rho\eta^2}{T^q}
% \] where \(q \leq 2l\), \(0< \eta \leq \frac{1}{2}\) and \(1 \leq \rho \leq \frac{1}{8\eta^2}\) are given constants that are independent of \(T\).
% Then, the upper bound of \eqref{Eq: the convergence rate of algorithm VR} is
% \[
% \text{Rate} = \calO (\max \{ T^{l-1}, T^{q-p-1}, T^{-2q}\})
% \]
% We can choose \(q = 2l = \frac{2}{3}\) and \(p=\frac{1}{3}\), then the upper bound is of the order \(\calO(T^{-\frac{2}{3})}\)
% \subsection{Discussion}
